\documentclass[a4paper]{article}

\usepackage[spanish]{babel}
\usepackage{graphicx}
\usepackage{amsmath, amssymb}
\usepackage[margin=2cm]{geometry}
\usepackage{fancyhdr}
\usepackage{enumerate}
\usepackage[shortlabels]{enumitem}
\usepackage{parskip}
\usepackage[most]{tcolorbox}
\usepackage{float}
\usepackage{booktabs}
\usepackage{mathrsfs}
\usepackage{tikz}
\usepackage[hidelinks]{hyperref}

\usetikzlibrary{calc}
\pagestyle{fancy}
\fancyhead[l]{Mecánica de Medios Continuos}
\fancyhead[c]{Taller 6}
\fancyhead[r]{\today}
\fancyfoot[c]{\thepage}
\setlength{\headheight}{14.5pt}
\renewcommand{\headrulewidth}{0.2pt}

\begin{document}
\tableofcontents

\section{Problema 1}

Considere un cilindro elástico de paredes gruesas, con radio interno $a$ y radio externo $b$, sometido al campo de esfuerzos:
$$ \sigma_{rr} = A + \frac{B}{r^2}, \qquad \sigma_{\theta\theta} = A - \frac{B}{r^2} $$
donde $A$ y $B$ son constantes y las demás componentes del tensor de esfuerzos son nulas.
\begin{enumerate}[(a)]
    \item Verifique que el campo de esfuerzos satisface la ecuación de equilibrio en ausencia de fuerzas volumétricas.
    \item Encuentre el vector de tracción en $r=a$, indicando la normal usada.
    \item Determine $A$ y $B$ si la superficie interna está sujeta a una presión uniforme $p$, mientras que la superficie externa está libre de tracción.
\end{enumerate}

\begin{figure}[H]
    \centering
    \definecolor{shellblue}{RGB}{80,145,195}
    \definecolor{pressurered}{RGB}{185,55,45}
    \begin{tikzpicture}[x=1cm,y=1cm,line cap=round,line join=round]
        \coordinate (O) at (3.1,2.4);
        \fill[shellblue!16] (O) circle (2.0);
        \fill[white] (O) circle (0.9);
        \draw[line width=1pt] (O) circle (2.0);
        \draw[line width=1pt] (O) circle (0.9);
        \fill (O) circle (0.035);
        \node[above left] at (O) {$z$};

        \draw[-latex,line width=0.8pt] (O) -- ++(35:0.9);
        \node[above right] at (3.65,2.95) {$a$};
        \draw[-latex,line width=0.8pt] (O) -- ++(-38:2.0);
        \node[below right] at (4.75,1.18) {$b$};

        \foreach \angle in {35,115,195,275} {
            \draw[pressurered,-latex,line width=0.9pt]
                ($(O)+(\angle:0.88)$) -- ($(O)+(\angle:1.30)$);
        }
        \node[pressurered,align=left] at (6.1,3.35) {presión interna\\uniforme $p$};
        \draw[pressurered,line width=0.6pt] (5.25,3.2) -- (4.1,2.95);
        \node[align=left] at (6.2,1.35) {pared externa\\sin tracción};
        \draw[black!55,line width=0.6pt] (5.3,1.55) -- (4.85,1.15);
        \node[align=center] at (8.0,2.4) {sección transversal\\del cilindro hueco};
    \end{tikzpicture}
    \caption{Geometría del cilindro de paredes gruesas y presión en su superficie interna.}
\end{figure}

\subsection{Solución}

\subsubsection*{Inciso (a)}
El equilibrio local, en ausencia de fuerzas volumétricas, se escribe como $\nabla\cdot\boldsymbol{\sigma}=\vec{0}$. Para ver su componente radial en coordenadas cilíndricas, usamos la expresión de la divergencia de un tensor en la base ortonormal $(\hat e_r,\hat e_\theta,\hat e_z)$:
$$
(\nabla\cdot\boldsymbol{\sigma})_r
=\frac{1}{r}\frac{\partial(r\sigma_{rr})}{\partial r}
+\frac{1}{r}\frac{\partial\sigma_{r\theta}}{\partial\theta}
+\frac{\partial\sigma_{rz}}{\partial z}
-\frac{\sigma_{\theta\theta}}{r}.
$$
El término $-\sigma_{\theta\theta}/r$ aparece porque la dirección $\hat e_\theta$ cambia con el ángulo. En este problema el campo es axisimétrico, no depende de $\theta$ ni de $z$, y las componentes cortantes son nulas. Entonces, $(\nabla\cdot\boldsymbol{\sigma})_r=0$ se reduce a
$$ \frac{\partial \sigma_{rr}}{\partial r} + \frac{\sigma_{rr} - \sigma_{\theta\theta}}{r} = 0 $$
Con el campo propuesto, calculamos
$$ \frac{\partial \sigma_{rr}}{\partial r} = -\frac{2B}{r^3} $$
$$ \sigma_{rr} - \sigma_{\theta\theta} = \left(A + \frac{B}{r^2}\right) - \left(A - \frac{B}{r^2}\right) = \frac{2B}{r^2} $$
Al sustituir ambos resultados en el equilibrio radial,
$$ -\frac{2B}{r^3} + \frac{1}{r}\left(\frac{2B}{r^2}\right) = -\frac{2B}{r^3} + \frac{2B}{r^3} = 0 $$
se verifica que la componente radial de la ecuación de equilibrio se satisface. Las componentes angular y axial también se anulan, pues no hay dependencias angulares o axiales ni esfuerzos cortantes.

\subsubsection*{Inciso (b)}
El vector de tracción $\vec{T}$ en una superficie con normal $\hat{n}$ se define como $\vec{T} = \boldsymbol{\sigma} \cdot \hat{n}$. Para la superficie $r=a$, si se toma $\hat{n}=\hat{e}_r$:
$$ \vec{T} = \sigma_{rr}(a)\hat{e}_r + \sigma_{r\theta}(a)\hat{e}_\theta + \sigma_{rz}(a)\hat{e}_z $$
Dado que las componentes de corte son nulas ($\sigma_{r\theta} = \sigma_{rz} = 0$), el vector de tracción es puramente radial:
$$ \vec{T} = \left(A + \frac{B}{a^2}\right) \hat{e}_r $$
Si $r=a$ es la pared interna, la normal exterior al material es $-\hat{e}_r$ y la tracción cambia de signo.

\subsubsection*{Inciso (c)}
Las condiciones de frontera son:
\begin{enumerate}[(i)]
    \item En la superficie interna ($r=a$), la presión uniforme impone $\sigma_{rr}(a)=-p$.
    \item En la superficie externa ($r=b$), la condición sin tracción impone $\sigma_{rr}(b)=0$.
\end{enumerate}

Esto genera el siguiente sistema de ecuaciones:
$$ A + \frac{B}{a^2} = -p $$
$$ A + \frac{B}{b^2} = 0 \implies A = -\frac{B}{b^2} $$
Sustituyendo $A$ en la primera ecuación:
$$ -\frac{B}{b^2} + \frac{B}{a^2} = -p \implies B \left(\frac{b^2 - a^2}{a^2 b^2}\right) = -p $$
Despejando las constantes $A$ y $B$:
$$ B = -\frac{p a^2 b^2}{b^2 - a^2} $$
$$ A = \frac{p a^2}{b^2 - a^2} $$

\section{Problema 2}

Justifique matemáticamente por qué una traslación y una rotación rígida no son deformaciones.
\begin{figure}[H]
    \centering
    \definecolor{moveblue}{RGB}{45,115,180}
    \definecolor{initialgray}{RGB}{115,120,125}
    \begin{tikzpicture}[x=1cm,y=1cm,line cap=round,line join=round]
        \node[font=\bfseries] at (2.5,3.1) {Traslación rígida};
        \draw[initialgray,dashed,line width=0.9pt] (0.6,0.8) rectangle (2.1,2.3);
        \fill[moveblue!14] (2.3,1.3) rectangle (3.8,2.8);
        \draw[moveblue,line width=1pt] (2.3,1.3) rectangle (3.8,2.8);
        \draw[moveblue,-latex,line width=0.9pt] (1.25,1.55) -- (2.85,2.05);
        \node[below] at (1.2,0.65) {inicial};
        \node[below,text=moveblue] at (3.0,1.15) {trasladada};

        \draw[black!22,line width=0.7pt] (5.0,0.45) -- (5.0,3.35);

        \node[font=\bfseries] at (7.5,3.35) {Rotación rígida};
        \draw[initialgray,dashed,line width=0.9pt] (6.2,1.1) rectangle (8.8,2.4);
        \begin{scope}[shift={(7.5,1.75)},rotate=28]
            \fill[moveblue!14] (-1.3,-0.65) rectangle (1.3,0.65);
            \draw[moveblue,line width=1pt] (-1.3,-0.65) rectangle (1.3,0.65);
        \end{scope}
        \draw[moveblue,-latex,line width=0.9pt] (8.4,2.8) arc[start angle=28,end angle=62,radius=0.8];
        \node[below] at (6.4,0.25) {inicial};
        \node[below,text=moveblue] at (8.6,0.25) {rotada};
    \end{tikzpicture}
    \caption{La traslación y la rotación cambian posición u orientación, pero no deforman el cuerpo.}
\end{figure}

\subsection{Solución}
En la teoría de pequeñas deformaciones, mostramos que $U_{ij} = \frac{1}{2}(\partial_j u_i + \partial_i u_j)$ se anula para una traslación y una rotación infinitesimal. Una rotación rígida finita también conserva exactamente longitudes y ángulos.

\subsubsection*{Traslación}
El campo de desplazamientos es un vector constante $\vec{u} = \vec{c}$, es decir, $u_i = c_i$.
$$ \partial_j u_i = \partial_j c_i = 0 $$
$$ U_{ij} = \frac{1}{2}(0 + 0) = 0 $$
Al ser el tensor de deformación nulo, no hay cambio en las longitudes ni en los ángulos; por lo tanto, no hay deformación.

\subsubsection*{Rotación}
Para una rotación infinitesimal con vector $\vec{\omega}$, el desplazamiento es $\vec{u} = \vec{\omega} \times \vec{x}$. En notación de índices: $u_i = \epsilon_{ijk} \omega_j x_k$.
Una rotación rígida finita se representa como $\vec{x}'=Q\vec{x}+\vec{c}$, donde $Q^{T}Q=I$. Para cualquier elemento de línea:
$$
\lVert d\vec{x}'\rVert^2 = (Q\,d\vec{x})\cdot(Q\,d\vec{x}) = d\vec{x}\cdot(Q^TQ)\,d\vec{x} = \lVert d\vec{x}\rVert^2.
$$
La derivada que se necesita aquí es $\partial_m u_i=\partial u_i/\partial x_m$ (la derivada de la componente de desplazamiento $u_i$ respecto a la coordenada $x_m$). En índices,
$$
\partial_m u_i
=\frac{\partial}{\partial x_m}\left(\epsilon_{ijk}\omega_jx_k\right)
=\epsilon_{ijk}\omega_j\frac{\partial x_k}{\partial x_m}
=\epsilon_{ijk}\omega_j\delta_{km}
=\epsilon_{ijm}\omega_j.
$$
Por ejemplo, escribiendo explícitamente $\vec{u}=\vec{\omega}\times\vec{x}$,
$$
\begin{aligned}
u_1&=\omega_2x_3-\omega_3x_2,\\
u_2&=\omega_3x_1-\omega_1x_3,\\
u_3&=\omega_1x_2-\omega_2x_1,
\end{aligned}
\qquad
\left[\partial_m u_i\right]
=\begin{pmatrix}
0&-\omega_3&\omega_2\\
\omega_3&0&-\omega_1\\
-\omega_2&\omega_1&0
\end{pmatrix}.
$$
Esta matriz es antisimétrica. Por eso, al tomar su parte simétrica para formar el tensor de deformación,
$$
U_{im}=\frac12(\partial_m u_i+\partial_i u_m)=0.
$$
Al anularse la parte simétrica del gradiente de desplazamientos, no existe deformación del material, solo un cambio de posición y orientación en el espacio.

\section{Problema 3}

A un cubo unitario cuyas aristas son paralelas a los ejes coordenados se le aplica el campo de desplazamientos $u_1 = kx_1$, $u_2 = u_3 = 0$ con $k = 10^{-4}$. Encuentre el incremento en la longitud de la diagonal en la cara paralela al plano $x_1x_2$.
\begin{figure}[H]
    \centering
    \definecolor{deformblue}{RGB}{40,115,185}
    \definecolor{initialgray}{RGB}{115,120,125}
    \begin{tikzpicture}[x=3.0cm,y=2.8cm,line cap=round,line join=round]
        \fill[deformblue!10] (0,0) -- (1.42,0) -- (1.42,1) -- (0,1) -- cycle;
        \draw[initialgray,dashed,line width=0.8pt] (0,0) rectangle (1,1);
        \draw[initialgray,dashed,line width=0.8pt] (0,0) -- (1,1);
        \draw[deformblue,line width=1pt] (0,0) rectangle (1.42,1);
        \draw[deformblue,-latex,line width=1.1pt] (0,0) -- (1.42,1);
        \draw[-latex,line width=0.8pt] (-0.08,0) -- (1.62,0) node[right] {$x_1$};
        \draw[-latex,line width=0.8pt] (0,-0.08) -- (0,1.22) node[above] {$x_2$};
        \node[initialgray,above left] at (0.42,0.43) {diagonal inicial};
        \node[deformblue,anchor=west] at (1.48,0.88) {diagonal deformada};
        \node[align=center] at (0.72,-0.38) {La deformación se exagera para\\facilitar la visualización.};
    \end{tikzpicture}
    \caption{La cara $x_1x_2$ se alarga en la dirección $x_1$, aumentando su diagonal.}
\end{figure}

\subsection{Solución}
El campo de desplazamientos es $u_1 = k x_1$, $u_2 = 0$, $u_3 = 0$.
El tensor de pequeñas deformaciones tiene una única componente no nula:
$$ U_{11} = \partial_1 u_1 = k $$
La diagonal de la cara paralela al plano $x_1 x_2$ de un cubo unitario conecta, por ejemplo, el origen $(0,0,0)$ con el punto $(1,1,0)$. El vector original es $\vec{L}_0 = \hat{e}_1 + \hat{e}_2$, y su longitud original es $L_0 = \sqrt{1^2 + 1^2} = \sqrt{2}$.
La diagonal deformada tiene vector $\vec{L}'=(1+k)\hat{e}_1+\hat{e}_2$, por lo que su longitud y el incremento exacto son
$$ L'=\sqrt{(1+k)^2+1}, \qquad \Delta L=\sqrt{(1+k)^2+1}-\sqrt{2}. $$
Para expandir a primer orden en $k$, factorizamos $2$ dentro de la raíz:
$$
L'=\sqrt{2+2k+k^2}
=\sqrt{2}\sqrt{1+k+\frac{k^2}{2}}.
$$
Usamos la expansión de Taylor $\sqrt{1+\varepsilon}=1+\varepsilon/2+O(\varepsilon^2)$ cuando $|\varepsilon|\ll1$. Con $\varepsilon=k+k^2/2$, a primer orden en $k$ queda
$$
L'=\sqrt{2}\left(1+\frac{k}{2}+O(k^2)\right),
\qquad
\Delta L=L'-L_0=\frac{k}{\sqrt{2}}+O(k^2).
$$
Para $k=10^{-4}$,
$$ \Delta L\approx\frac{10^{-4}}{\sqrt{2}}\approx 7.07\times10^{-5}. $$

\section{Problema 4}

Considere el campo de desplazamientos $u_1 = k(2x_1^2 + x_1x_2)$, $u_2 = kx_2^2$, $u_3 = 0$ con $k = 10^{-4}$.
\begin{enumerate}[(a)]
    \item Encuentre las elongaciones unitarias y el cambio en el ángulo para dos elementos de línea materiales $d\vec{x}_1 = dx_1\hat{e}_1$ y $d\vec{x}_2 = dx_2\hat{e}_2$ para la partícula en $\vec{x} = \hat{e}_1 - \hat{e}_2$.
    \item Bosqueje las posiciones deformadas de estos dos elementos.
\end{enumerate}

\subsection{Solución}

\subsubsection*{Inciso (a)}
El campo de desplazamientos es $u_1 = k(2x_1^2 + x_1 x_2)$, $u_2 = k x_2^2$, $u_3 = 0$.
Calculamos las derivadas parciales del desplazamiento:
$$ \partial_1 u_1 = k(4x_1 + x_2), \quad \partial_2 u_1 = k x_1, \quad \partial_3 u_1 = 0 $$
$$ \partial_1 u_2 = 0, \quad \partial_2 u_2 = 2k x_2, \quad \partial_3 u_2 = 0 $$
Las componentes del tensor de deformación $U_{ij} = \frac{1}{2}(\partial_j u_i + \partial_i u_j)$ son:
$$ U_{11} = k(4x_1 + x_2), \quad U_{22} = 2k x_2, \quad U_{12} = \frac{1}{2}k x_1 $$
Evaluamos en la partícula $\vec{x} = \hat{e}_1 - \hat{e}_2$, lo que significa $x_1 = 1$ y $x_2 = -1$:
$$ U_{11} = k(4(1) - 1) = 3k $$
$$ U_{22} = 2k(-1) = -2k $$
$$ U_{12} = \frac{1}{2}k(1) = 0.5k $$

Las elongaciones unitarias para los elementos materiales en las direcciones de los ejes son las componentes diagonales del tensor de deformación:
\begin{enumerate}
    \item Para $d\vec{x}^{(1)} = dx_1 \hat{e}_1$: $\epsilon_1 = U_{11} = 3k = 3 \times 10^{-4}$.
    \item Para $d\vec{x}^{(2)} = dx_2 \hat{e}_2$: $\epsilon_2 = U_{22} = -2k = -2 \times 10^{-4}$.
\end{enumerate}

La deformación por corte ingenieril es
$$ \gamma_{12} = 2 U_{12} = k = 10^{-4}. $$
Para construir los vectores que determinan el ángulo, usamos el gradiente del desplazamiento evaluado en el punto $(x_1,x_2)=(1,-1)$:
$$
\nabla\vec u=
\begin{pmatrix}
\partial_1u_1&\partial_2u_1\\
\partial_1u_2&\partial_2u_2
\end{pmatrix}_{(1,-1)}
=\begin{pmatrix}3k&k\\0&-2k\end{pmatrix}.
$$
Un elemento diferencial se transforma según $d\vec{x}'=(I+\nabla\vec{u})d\vec{x}$. Para los elementos iniciales $d\vec{x}_1=dx_1\hat e_1$ y $d\vec{x}_2=dx_2\hat e_2$, dividimos los vectores deformados por sus longitudes iniciales diferenciales; así obtenemos los vectores tangentes por unidad de longitud original:
$$
\begin{aligned}
\vec v_1
&=\frac{d\vec{x}'_1}{dx_1}
=(I+\nabla\vec u)\hat e_1
=\begin{pmatrix}1+3k&k\\0&1-2k\end{pmatrix}
\begin{pmatrix}1\\0\end{pmatrix}
=(1+3k)\hat e_1,\\[2mm]
\vec v_2
&=\frac{d\vec{x}'_2}{dx_2}
=(I+\nabla\vec u)\hat e_2
=\begin{pmatrix}1+3k&k\\0&1-2k\end{pmatrix}
\begin{pmatrix}0\\1\end{pmatrix}
=k\hat e_1+(1-2k)\hat e_2.
\end{aligned}
$$
No son vectores unitarios deformados; sus direcciones sí son las de los elementos después de deformarse, y por eso sirven para hallar el ángulo.
El primero sigue sobre el eje $x_1$. El segundo forma con ese eje el ángulo
$$
\theta'=\arctan\!\left(\frac{1-2k}{k}\right)
=\frac{\pi}{2}-\arctan\!\left(\frac{k}{1-2k}\right),
$$
para el valor positivo y pequeño de $k$ usado aquí. Como el ángulo inicial es $\theta_0=\pi/2$, el cambio es
$$
\Delta\theta=\theta'-\theta_0
=-\arctan\!\left(\frac{k}{1-2k}\right)
\approx -k=-10^{-4}\,\text{rad}.
$$
El signo negativo indica que el ángulo recto se cierra. A primer orden, esto coincide con $\Delta\theta\approx-\gamma_{12}$.

\subsubsection*{Inciso (b)}
Para describir las posiciones deformadas, usamos la relación $d\vec{x}' = d\vec{x} + (\nabla \vec{u}) d\vec{x}$. El gradiente de desplazamiento evaluado en el punto es:
$$ \nabla \vec{u} = \begin{pmatrix} 3k & k & 0 \\ 0 & -2k & 0 \\ 0 & 0 & 0 \end{pmatrix} $$
Si los extremos del elemento están en $\vec{x}$ y $\vec{x}+d\vec{x}$, el cambio de desplazamiento entre ellos es
$$
\Delta\vec{u}=\vec{u}(\vec{x}+d\vec{x})-\vec{u}(\vec{x})
\approx(\nabla\vec{u})\,d\vec{x},
\qquad
\Delta u_i=\frac{\partial u_i}{\partial x_j}\,dx_j.
$$
Esta es la expansión de Taylor a primer orden en el tamaño diferencial $d\vec{x}$. Para el primer elemento, $d\vec{x}_1=dx_1\hat e_1$, así que
$$
\Delta\vec{u}_1
=\begin{pmatrix}3k&k&0\\0&-2k&0\\0&0&0\end{pmatrix}
\begin{pmatrix}dx_1\\0\\0\end{pmatrix}
=\begin{pmatrix}3k\,dx_1\\0\\0\end{pmatrix}
=3k\,dx_1\hat e_1,
\qquad
d\vec{x}'_1=d\vec{x}_1+\Delta\vec{u}_1=(1+3k)dx_1\hat e_1.
$$
Para el segundo elemento, $d\vec{x}_2=dx_2\hat e_2$, de modo que su componente en $x_1$ es cero y la multiplicación da
$$
\Delta\vec{u}_2
=\begin{pmatrix}3k&k&0\\0&-2k&0\\0&0&0\end{pmatrix}
\begin{pmatrix}0\\dx_2\\0\end{pmatrix}
=\begin{pmatrix}k\,dx_2\\-2k\,dx_2\\0\end{pmatrix}
=dx_2(k\hat e_1-2k\hat e_2).
$$
Por tanto,
$$
d\vec{x}'_2=d\vec{x}_2+\Delta\vec{u}_2
=dx_2\left[k\hat e_1+(1-2k)\hat e_2\right].
$$
La figura muestra las configuraciones inicial y deformada; la inclinación y los cambios de longitud se exageran para hacerlos visibles.
\begin{figure}[H]
    \centering
    \definecolor{deformblue}{RGB}{40,115,185}
    \definecolor{initialgray}{RGB}{115,120,125}
    \begin{tikzpicture}[x=3.5cm,y=2.7cm,line cap=round,line join=round]
        \coordinate (O) at (0,0);
        \draw[-latex,line width=0.8pt] (-0.2,0) -- (1.85,0) node[right] {$x_1$};
        \draw[-latex,line width=0.8pt] (0,-0.2) -- (0,1.55) node[above] {$x_2$};

        \draw[initialgray,dashed,line width=1pt,-latex] (O) -- (1.05,0)
            node[below right] {$d\vec{x}_1$};
        \draw[initialgray,dashed,line width=1pt,-latex] (O) -- (0,1.05)
            node[above left] {$d\vec{x}_2$};

        \draw[deformblue,line width=1.3pt,-latex] (O) -- (1.55,0)
            node[below] {$d\vec{x}'_1$};
        \draw[deformblue,line width=1.3pt,-latex] (O) -- (0.38,0.92)
            node[above right] {$d\vec{x}'_2$};

        \draw[initialgray,line width=0.6pt] (0.25,0) arc[start angle=0,end angle=90,radius=0.25];
        \node[initialgray] at (0.4,0.38) {$90^\circ$};
        \draw[deformblue,line width=0.7pt] (0.45,0) arc[start angle=0,end angle=68,radius=0.45];
        \node[deformblue] at (0.72,0.48) {$\theta'<90^\circ$};

        \node[initialgray,align=left] at (2.45,1.0) {inicial\\(discontinua)};
        \node[deformblue,align=left] at (2.45,0.45) {deformado\\(continua)};
    \end{tikzpicture}
    \caption{Elementos materiales en el punto considerado, antes y después de la deformación.}
\end{figure}

\section{Problema 5}

Dadas las componentes no nulas del tensor de deformación $U_{11} = \frac{1}{\alpha} f(x_2, x_3)$, $U_{22} = U_{33} = -\frac{\nu}{\alpha} f(x_2, x_3)$, y $U_{12} = U_{13} = U_{23} = 0$, demuestre que las condiciones de compatibilidad requieren que $f(x_2, x_3)$ sea lineal en $x_2$ y $x_3$.

\subsection{Solución}
Suponemos $\alpha\ne 0$, pues aparece en el denominador de las deformaciones. Las ecuaciones de compatibilidad de Saint-Venant son las condiciones que debe satisfacer $U_{ij}$ para que exista un campo de desplazamientos $u_i$ tal que $U_{ij}=\tfrac12(\partial_j u_i+\partial_i u_j)$. En notación compacta se escriben
$$
U_{ij,kl}+U_{kl,ij}-U_{ik,jl}-U_{jl,ik}=0,
\qquad i,j,k,l\in\{1,2,3\},
$$
donde la coma indica derivación: por ejemplo, $U_{11,22}=\partial^2 U_{11}/\partial x_2^2$. En tres dimensiones, las seis ecuaciones independientes pueden escribirse como
$$
\begin{aligned}
U_{11,22}+U_{22,11}&=2U_{12,12},\\
U_{22,33}+U_{33,22}&=2U_{23,23},\\
U_{33,11}+U_{11,33}&=2U_{13,13},\\
U_{12,33}+U_{33,12}&=U_{13,23}+U_{23,13},\\
U_{23,11}+U_{11,23}&=U_{12,13}+U_{13,12},\\
U_{13,22}+U_{22,13}&=U_{12,23}+U_{23,12}.
\end{aligned}
$$

Ahora aplicamos este sistema al campo dado:
$$
U_{11}=\frac{f(x_2,x_3)}{\alpha},\qquad
U_{22}=U_{33}=-\frac{\nu}{\alpha}f(x_2,x_3),\qquad
U_{12}=U_{13}=U_{23}=0.
$$
Como $f$ depende solo de $x_2$ y $x_3$, cualquier derivada respecto a $x_1$ es cero. Además, las componentes de corte y sus derivadas también son cero. Por eso:
\begin{itemize}
    \item La primera ecuación queda $U_{11,22}=0$, es decir, $f_{,22}/\alpha=0$.
    \item La tercera ecuación queda $U_{11,33}=0$, es decir, $f_{,33}/\alpha=0$.
    \item La quinta ecuación queda $U_{11,23}=0$, es decir, $f_{,23}/\alpha=0$.
\end{itemize}
Como $\alpha\ne0$, las condiciones que resultan son
$$
f_{,22}=0,\qquad f_{,33}=0,\qquad f_{,23}=0,
$$
o, de forma explícita,
$$
\frac{\partial^2 f}{\partial x_2^2}=0,\qquad
\frac{\partial^2 f}{\partial x_3^2}=0,\qquad
\frac{\partial^2 f}{\partial x_2\partial x_3}=0.
$$
Las otras tres ecuaciones no añaden restricciones: la segunda se satisface al usar $f_{,22}=f_{,33}=0$, y la cuarta y la sexta contienen derivadas respecto a $x_1$ o derivadas de componentes de corte nulas.

Para ver qué forma tiene $f$, integramos estas condiciones. De $f_{,22}=0$ se sigue que $f$ debe ser lineal en $x_2$, aunque sus coeficientes podrían depender de $x_3$:
$$
f(x_2,x_3)=A(x_3)x_2+B(x_3).
$$
La condición $f_{,33}=0$ implica $A''(x_3)x_2+B''(x_3)=0$, de modo que $A''=B''=0$. La condición mixta $f_{,23}=0$ exige además $A'(x_3)=0$, por lo que $A$ es constante y $B$ es lineal en $x_3$. En consecuencia,
$$ f(x_2, x_3) = A x_2 + B x_3 + C $$
con $A$, $B$ y $C$ constantes. Así, las condiciones de compatibilidad obligan a que $f$ sea afín (lineal más una constante) en $x_2$ y $x_3$.

\section{Problema 6}

Demuestre que para un material isotrópico que sigue la ley de Hooke las direcciones principales de esfuerzo coinciden con las direcciones principales de deformación, y encuentre una relación entre los valores de los esfuerzos principales y las deformaciones principales.

\subsection{Solución}
La ley de Hooke para un material isotrópico linealmente elástico es
$$
\boldsymbol{\sigma}
=2\mu\mathbf{U}+\lambda\,\operatorname{tr}(\mathbf{U})\mathbf{I},
$$
donde $\mathbf{U}$ es el tensor de pequeñas deformaciones, $\boldsymbol{\sigma}$ es el tensor de esfuerzos, $\mathbf{I}$ es el tensor identidad y $\lambda,\mu$ son los coeficientes de Lamé. En componentes, $\operatorname{tr}(\mathbf{U})=U_{kk}$.

El tensor $\mathbf{U}$ es simétrico, así que tiene tres direcciones principales ortogonales. Llamemos $\hat{n}_i$ a una de ellas y $U_i$ a la deformación principal correspondiente. Decir que son un par propio significa, por definición, que al aplicar el tensor a ese vector, el resultado queda en la misma dirección:
$$
\mathbf{U}\hat{n}_i=U_i\hat{n}_i.
$$
Aquí $\hat{n}_i$ indica la dirección, mientras que el escalar $U_i$ indica cuánto se estira o contrae esa dirección (y con qué signo). Esta propiedad es justamente lo que define un vector propio y su valor propio. Las direcciones principales de deformación son, por tanto, los vectores propios de $\mathbf{U}$.

Ahora aplicamos el tensor de esfuerzos a esa misma dirección. Usamos que la identidad no cambia ningún vector, $\mathbf{I}\hat{n}_i=\hat{n}_i$, y sustituimos el par propio de deformación:
$$
\begin{aligned}
\boldsymbol{\sigma}\hat{n}_i
&=\left[2\mu\mathbf{U}+\lambda\,\operatorname{tr}(\mathbf{U})\mathbf{I}\right]\hat{n}_i\\
&=2\mu\left(\mathbf{U}\hat{n}_i\right)
 +\lambda\,\operatorname{tr}(\mathbf{U})\left(\mathbf{I}\hat{n}_i\right)\\
&=2\mu U_i\hat{n}_i
 +\lambda\,\operatorname{tr}(\mathbf{U})\hat{n}_i\\
&=\left[2\mu U_i+\lambda\,\operatorname{tr}(\mathbf{U})\right]\hat{n}_i.
\end{aligned}
$$
El resultado vuelve a ser un múltiplo escalar de $\hat{n}_i$. Por definición, esto demuestra que $\hat{n}_i$ también es un vector propio de $\boldsymbol{\sigma}$: la dirección no cambia al aplicar el tensor de esfuerzos. El término $2\mu\mathbf U$ conserva las direcciones propias de $\mathbf U$, y el término isotrópico $\lambda\operatorname{tr}(\mathbf U)\mathbf I$ solo añade el mismo escalar a todas las direcciones, sin privilegiar ni girar ninguna.

Como $\mathbf U$ es simétrico, sus vectores propios pueden elegirse ortonormales y su traza es la suma de sus valores propios:
$$
\operatorname{tr}(\mathbf U)=U_1+U_2+U_3.
$$
Por consiguiente, el valor propio de esfuerzo en la dirección $\hat n_i$ es
$$
\boxed{\sigma_i=2\mu U_i+\lambda(U_1+U_2+U_3)},
\qquad i=1,2,3.
$$
En la base principal común, esto equivale a
$$
\mathbf U=\begin{pmatrix}U_1&0&0\\0&U_2&0\\0&0&U_3\end{pmatrix},
\qquad
\boldsymbol{\sigma}=\begin{pmatrix}\sigma_1&0&0\\0&\sigma_2&0\\0&0&\sigma_3\end{pmatrix}.
$$
Así, las direcciones principales de esfuerzo y deformación coinciden. Si algún valor propio está repetido, las direcciones dentro de ese subespacio no son únicas, pero se puede elegir una misma base ortonormal principal para ambos tensores.

La figura ilustra en dos dimensiones esta acción: cada tensor puede cambiar la magnitud del vector asociado a una dirección principal, pero su imagen permanece sobre esa misma dirección (si el valor propio es negativo, el sentido de la flecha se invierte, no el eje principal).
\begin{figure}[H]
    \centering
    \definecolor{strainblue}{RGB}{40,115,185}
    \definecolor{stressred}{RGB}{185,55,45}
    \definecolor{principalgray}{RGB}{125,130,135}
    \begin{tikzpicture}[x=1cm,y=1cm,line cap=round,line join=round]
        % Paneles
        \draw[black!35,rounded corners=4pt] (0,0) rectangle (5.15,3.25);
        \draw[black!35,rounded corners=4pt] (5.45,0) rectangle (10.60,3.25);
        \node[font=\bfseries] at (2.575,2.88) {Tensor de deformación $\mathbf{U}$};
        \node[font=\bfseries] at (8.025,2.88) {Tensor de esfuerzo $\boldsymbol{\sigma}$};

        % Direcciones principales de deformación y sus imágenes bajo U
        \coordinate (OU) at (1.0,0.85);
        \draw[principalgray,dashed,-latex,line width=0.8pt] (OU) -- (4.55,0.85);
        \draw[principalgray,dashed,-latex,line width=0.8pt] (OU) -- (1.0,2.35);
        \draw[strainblue,-latex,line width=1.5pt] (1.0,1.00) -- (2.85,1.00);
        \draw[strainblue,-latex,line width=1.5pt] (1.15,0.85) -- (1.15,1.95);
        \node[principalgray,right] at (4.55,0.85) {$\hat{n}_1$};
        \node[principalgray,above] at (1.0,2.35) {$\hat{n}_2$};
        \node[strainblue,below] at (2.25,0.83) {imagen};
        \fill (OU) circle (0.035);

        % Las mismas direcciones principales y sus imágenes bajo sigma
        \coordinate (OS) at (6.45,0.85);
        \draw[principalgray,dashed,-latex,line width=0.8pt] (OS) -- (10.0,0.85);
        \draw[principalgray,dashed,-latex,line width=0.8pt] (OS) -- (6.45,2.35);
        \draw[stressred,-latex,line width=1.5pt] (6.45,1.00) -- (8.75,1.00);
        \draw[stressred,-latex,line width=1.5pt] (6.60,0.85) -- (6.60,2.10);
        \node[principalgray,right] at (10.0,0.85) {$\hat{n}_1$};
        \node[principalgray,above] at (6.45,2.35) {$\hat{n}_2$};
        \node[stressred,below] at (7.60,0.83) {imagen};
        \fill (OS) circle (0.035);
    \end{tikzpicture}
    \caption{Ambos tensores envían cada dirección principal a un vector paralelo a ella. Las longitudes de las flechas son esquemáticas y se muestran valores propios positivos; un valor propio negativo invertiría el sentido, pero no la dirección principal.}
\end{figure}

\section{Problema 7}

Un material es incompresible si su volumen no cambia bajo ningún estado de esfuerzo. Pruebe que para un sólido incompresible, isotrópico y linealmente elástico con módulo de Young finito $E_Y$: (a) la razón de Poisson es $\nu = 1/2$, (b) el módulo de corte es $\mu = E_Y/3$ y (c) el módulo volumétrico $K\to\infty$ y $\lambda\to\infty$, mientras $K-\lambda=2\mu/3$.

\subsection{Solución}

\subsubsection*{Inciso (a): Razón de Poisson}
La deformación volumétrica (dilatación) $e$ se define como la traza del tensor de deformaciones, $e = U_{kk}$. Usando la ley de Hooke invertida, la dilatación en función de la traza del tensor de esfuerzos $\sigma_{kk}$ es:
$$ e = U_{kk} = \frac{1 - 2\nu}{E_Y} \sigma_{kk} $$
Para un material incompresible, el volumen no cambia bajo ningún estado de esfuerzo, lo que significa que $e = 0$ para cualquier valor arbitrario de $\sigma_{kk}$. Dado que el módulo de Young $E_Y$ es finito y distinto de cero, la única forma de que esta ecuación se cumpla universalmente es que el numerador sea cero:
$$ 1 - 2\nu = 0 \implies \nu = \frac{1}{2} $$

\subsubsection*{Inciso (b): Módulo de corte}
La relación fundamental entre el módulo de Young $E_Y$, el módulo de corte $\mu$ y la razón de Poisson $\nu$ es:
$$ E_Y = 2\mu(1 + \nu) $$
Sustituyendo el valor $\nu = 1/2$ obtenido en el inciso anterior:
$$ E_Y = 2\mu \left(1 + \frac{1}{2}\right) = 2\mu \left(\frac{3}{2}\right) = 3\mu $$
Despejando el módulo de corte, obtenemos directamente:
$$ \mu = \frac{E_Y}{3} $$

\subsubsection*{Inciso (c): Comportamiento de los módulos $K$ y $\lambda$}
El módulo de compresibilidad volumétrica $K$ se relaciona con $E_Y$ y $\nu$ mediante:
$$ K = \frac{E_Y}{3(1 - 2\nu)} $$
Al tender $\nu \to 1/2$, el denominador tiende a cero mientras el numerador es finito y positivo, por lo que:
$$ K \to \infty $$
La relación entre $K$, el parámetro de Lamé $\lambda$ y el módulo de corte $\mu$ es:
$$ K = \lambda + \frac{2}{3}\mu $$
De aquí, $\lambda=K-2\mu/3\to\infty$, pues $\mu=E_Y/3$ permanece finito. Además,
$$ K - \lambda = \frac{2}{3}\mu $$
Esta diferencia es estrictamente finita e igual a $\frac{2}{3}\mu$, como se quería demostrar. Físicamente, aunque $K$ y $\lambda$ divergen en el límite incompresible, su diferencia permanece finita y está determinada por el módulo de corte $\mu$.

\section{Problema 8}

Encuentre la deformación de un sólido isotrópico elástico sumergido en un fluido a presión constante $p$.

\subsection{Solución}
Al estar sumergido en un fluido a presión $p$, el sólido experimenta un estado de esfuerzo hidrostático en todas sus direcciones. El tensor de esfuerzos es:
$$ \sigma_{ij} = -p \delta_{ij} $$
La traza del tensor de esfuerzos es $\sigma_{kk} = -3p$.
Utilizando la ley de Hooke generalizada para materiales isotrópicos, el tensor de deformación $U_{ij}$ se relaciona con los esfuerzos mediante:
$$ U_{ij} = \frac{1+\nu}{E_Y} \sigma_{ij} - \frac{\nu}{E_Y} \sigma_{kk} \delta_{ij} $$
Sustituyendo $\sigma_{ij}$:
$$ U_{ij} = \frac{1+\nu}{E_Y} (-p \delta_{ij}) - \frac{\nu}{E_Y} (-3p) \delta_{ij} = -\frac{p(1-2\nu)}{E_Y} \delta_{ij} $$
Definiendo el módulo de compresibilidad volumétrica como $K = \frac{E_Y}{3(1-2\nu)}$, la deformación se expresa de forma compacta como:
$$ U_{ij} = -\frac{p}{3K} \delta_{ij} $$
Esto indica que el sólido sufre una compresión uniforme. Las componentes diagonales del tensor de deformación son $U_{xx} = U_{yy} = U_{zz} = -\frac{p}{3K}$, y las componentes de corte son nulas.

Un campo de desplazamientos particular que produce esta deformación uniforme es:
$$ \vec{u} = -\frac{p}{3K}\vec{r}. $$
La solución general puede sumarle cualquier traslación y rotación infinitesimal de cuerpo rígido.

\section{Problema 9}

Encuentre la deformación de una varilla que es estirada al ejercer sobre sus extremos un esfuerzo $p$ en términos del módulo de Young y la razón de Poisson.

\subsection{Solución}
Alineamos el eje $z$ con el eje longitudinal de la varilla. Al aplicar un esfuerzo de tracción $p$ uniformemente distribuido en los extremos y no haber fuerzas laterales, el único componente no nulo del tensor de esfuerzos es:
$$ \sigma_{zz} = p, \quad \sigma_{ij} = 0 \text{ para } (i,j) \neq (z,z) $$
Aplicando la ley de Hooke invertida, las componentes del tensor de deformación son:
$$ U_{zz} = \frac{p}{E_Y} $$
$$ U_{xx} = U_{yy} = -\nu \frac{p}{E_Y} $$
$$ U_{xy} = U_{yz} = U_{zx} = 0 $$
Un campo de desplazamientos que reproduce estas deformaciones es
$$
\vec{u}= -\nu\frac{p}{E_Y}x\,\hat{e}_x
          -\nu\frac{p}{E_Y}y\,\hat{e}_y
          +\frac{p}{E_Y}z\,\hat{e}_z.
$$
Puede añadirse un desplazamiento de cuerpo rígido sin alterar la deformación.
Físicamente, esto describe una varilla que se alarga en la dirección $z$ y se contrae uniformemente en las direcciones transversales $x$ e $y$ debido al efecto de Poisson.

\section{Problema 10}

Determine la deformación de una varilla larga de longitud $l$ en un campo gravitacional constante.

\subsection{Solución}
Modelamos la varilla como esbelta, con esfuerzo uniaxial y sin considerar los efectos locales de los extremos. El eje $z$ apunta hacia arriba, desde $z=0$ hasta $z=l$, y la gravedad actúa en dirección $-z$ con aceleración $g_0$. La densidad del material es $\rho$.
La ecuación de equilibrio $\partial_j\sigma_{ij}+f_i=0$, con $f_z=-\rho g_0$, se reduce a
$$ \partial_z \sigma_{zz} = \rho g_0. $$
Integrando y aplicando la condición de frontera de superficie libre en el tope ($z=l$), donde $\sigma_{zz}(l) = 0$, obtenemos:
$$ \sigma_{zz} = -\rho g_0 (l - z) $$
Todos los demás componentes del tensor de esfuerzos son nulos ($\sigma_{ij} = 0$).
Las componentes del tensor de deformación se calculan mediante la ley de Hooke:
$$ U_{zz} = \frac{\sigma_{zz}}{E_Y} = -\frac{\rho g_0}{E_Y} (l - z) $$
$$ U_{xx} = U_{yy} = -\nu \frac{\sigma_{zz}}{E_Y} = \frac{\nu \rho g_0}{E_Y} (l - z) $$
$$ U_{xy} = U_{yz} = U_{zx} = 0 $$
Las deformaciones son compatibles con el siguiente campo de desplazamientos particular:
$$
\begin{aligned}
u_x&=\frac{\nu\rho g_0}{E_Y}(l-z)x,\\
u_y&=\frac{\nu\rho g_0}{E_Y}(l-z)y,\\
u_z&=-\frac{\rho g_0}{E_Y}\left(lz-\frac{z^2}{2}\right)
     +\frac{\nu\rho g_0}{2E_Y}(x^2+y^2).
\end{aligned}
$$
Se fija $\vec{u}(0,0,0)=\vec{0}$ como elección del movimiento rígido global; esto no equivale a imponer un empotramiento en toda la base. El término axial es negativo y representa el acortamiento por el peso propio; los términos transversales describen la expansión lateral de Poisson y hacen compatible el campo de deformaciones con los cortes nulos.

\end{document}